\documentclass[a4paper,12pt]{article}
\usepackage[utf8]{inputenc}
\usepackage[english]{babel}
\usepackage{authblk}
\usepackage{graphicx}
\usepackage{mathptmx}
\usepackage[singlespacing]{setspace}
\usepackage[headheight=1in,margin=1in]{geometry}
\usepackage{fancyhdr}


\renewcommand{\headrulewidth}{0pt}
\pagestyle{fancy}


\makeatletter
\def\@maketitle{%
  \newpage

  \begin{center}%
  \let \footnote \thanks
    {\LARGE \@title \par}%
  \end{center}%
  \par
  \vskip 0.1em}
\makeatother

\chead{%
  CS2Italy – May 19-21, 2026 – Torino, Italy\\
  %
}

\graphicspath{{images/}}

\title{From Static Snapshots to Living Databases: Automated Policy Classification for Dynamic Social Science Datasets}


\date{}

\begin{document}

\maketitle
\thispagestyle{fancy}

\begin{center}
\textit{Author names and affiliations to be added}
\end{center}

\section*{Introduction}

Social science research increasingly relies on large-scale policy databases to understand governance patterns and inform evidence-based decision-making. However, most policy databases face a fundamental challenge: they are static snapshots that quickly become outdated as new policies are enacted. The Food System Policy Database (FSPD) \cite{lowder2025}, containing approximately 11,900 policies across 192 countries, exemplifies both the value and limitation of such resources. While it provides unprecedented insight into global food system governance, maintaining its currency requires continuous manual curation and classification---a labor-intensive process that limits scalability and timeliness.

This challenge extends beyond food policy to any domain where social science research generates structured databases from administrative or policy documents. The conventional workflow---collect, clean, classify, publish, repeat---treats data curation as project-specific rather than as a generalizable methodological problem. Yet the classification task itself (assigning policies to categories, identifying objectives, extracting structured features from text) is fundamentally algorithmic and amenable to automation.

We present a computational approach that transforms the FSPD from a static resource into a ``living database'' through automated machine learning-based classification. More importantly, we demonstrate how the methodology itself---independent of the specific policy domain---can be adapted to any research context where experts have created structured classifications of text-based social data. Our contribution is thus twofold: an applied solution for food policy research, and a transferable framework for automating the maintenance of expert-curated social science databases.

\section*{Data and Methods}

The FSPD aggregates policies from four source databases (FAPDA, PINE, GIfNA, Nourishing) spanning 195 countries. Each policy includes a text description (\textit{policydecision\_details}, average 603 characters) and expert-assigned classifications across multiple dimensions: (1) \textit{segment} (11 categories: inputs, farming, trade, processing, retail, food service, final consumption, etc.); (2) \textit{lever} (113 specific policy instruments such as subsidies, taxes, regulations); (3) \textit{type} (9 high-level categories of policy intervention); and (4) three binary objectives (inclusion, environmental sustainability, dietary health).

The dataset presents several computational challenges characteristic of real-world social science data: severe class imbalance (lever classifications have a 246:1 ratio between most and least frequent classes), multilingual text (Spanish, French, English), and the need for simultaneous prediction across multiple interdependent taxonomies. These challenges mirror those faced in many policy and administrative text classification tasks across social science domains.

We developed a multi-task learning approach using XLM-RoBERTa-base \cite{conneau2020}, a multilingual transformer model (278M parameters) that handles text in over 100 languages without translation. The architecture features six prediction heads sharing a common text encoder: three multi-class classifiers (segment, lever, type) and three binary classifiers (policy objectives). This multi-task design enables the model to learn shared semantic representations while respecting the hierarchical and interdependent nature of policy classifications.

To address severe class imbalance, we implemented inverse frequency weighting for the three multi-class tasks, giving higher loss contribution to underrepresented policy instruments. Training employed standard transformer fine-tuning practices: AdamW optimization with weight decay (0.01), linear learning rate scheduling with warmup, gradient clipping (max norm 1.0), and early stopping based on validation loss. The model was trained on [XX]\% of the data, with [XX]\% validation and [XX]\% held-out test sets, stratified by policy segment to ensure representative splits.

Critically, our pipeline includes confidence-based quality control: each prediction generates a calibrated probability score, and predictions below a threshold (0.7) are automatically flagged for expert review. This human-in-the-loop design acknowledges that fully automated classification may miss nuanced cases while dramatically reducing the manual workload.

\section*{Results}

The trained model achieved strong performance across all six classification tasks on the held-out test set: segment accuracy [XX]\% (F1: [XX]), lever accuracy [XX]\% (F1: [XX]), type accuracy [XX]\% (F1: [XX]), and binary objective accuracies ranging from [XX--XX]\%. Weighted F1 scores account for class imbalance, with the lever classification---the most granular and imbalanced task---achieving [XX]\% weighted F1, indicating effective learning despite the 113-class problem with extreme frequency skew.

Confidence score analysis revealed that approximately [XX]\% of predictions exceeded the 0.7 threshold across all tasks, meaning the system can confidently classify the majority of new policies without human intervention. For low-confidence predictions, manual review maintains classification quality while requiring only a fraction of the effort needed for fully manual curation.

Cross-database analysis showed consistent performance across the source databases (FAPDA, PINE, GIfNA), suggesting the model generalizes well to different policy documentation styles and thematic focuses. Multilingual evaluation confirmed robust performance on Spanish and French texts, validating the choice of XLM-RoBERTa for international policy data.

\section*{Discussion and Implications}

This work demonstrates a practical pathway for transforming expert-curated social science databases from static snapshots into automatically maintained resources. The methodology is deliberately domain-agnostic: while applied here to food policy, the same pipeline---multilingual transformer + multi-task learning + imbalance handling + confidence-based review---can be adapted to legislative databases, administrative records, judicial decisions, news archives, or any corpus where human experts have established classification schemes.

The broader implication is that research projects producing structured social science data can build sustainability into their data collection processes. Rather than treating database construction as a one-time project deliverable, researchers can develop classification models alongside their initial curation efforts, ensuring that data remains current with minimal ongoing effort. This shifts computational social science from retroactive analysis of static data toward continuous monitoring of evolving social phenomena.

We envision extending this work in two directions. First, integrating categorical metadata (country, region, income level, temporal features) through a two-stage architecture: using transformer-derived text embeddings as input to gradient boosting models (XGBoost) that can natively handle mixed text and tabular features. Second, adapting the pipeline for cross-domain transfer learning, testing whether models trained on food policy classifications can accelerate training for related domains (environmental policy, agricultural regulation, public health interventions).

The living FSPD pipeline will be deployed with monthly automated updates ingesting new policies from source databases, with flagged predictions routed to domain experts. This demonstrates a sustainable model for maintaining research infrastructure beyond initial project funding. More importantly, by releasing the classification pipeline as open-source tooling, we provide computational social scientists with reusable infrastructure for automating their own data curation workflows, reducing barriers to longitudinal and real-time analysis of social and policy phenomena.

\section*{References}

\small
\begin{itemize}
\itemsep0pt
\item Conneau, A., Khandelwal, K., Goyal, N., et al. (2020). Unsupervised Cross-lingual Representation Learning at Scale. \textit{Proceedings of ACL}, 8440--8451.
\item Lowder, S.K., Ruggeri Laderchi, C., Cerutti, N., Parsons, K. (2025). Food System Policies: a Global Snapshot from the Food System Policy Database (FSPD). \textit{Global Food Security}. https://doi.org/10.1016/j.gfs.2025.100894
\end{itemize}

\newpage

\begin{figure}[htp]
\centering
\fbox{\parbox{0.9\textwidth}{\centering\vspace{100pt}\textbf{[FIGURE SUGGESTION]}\\
\textbf{Option 1:} Model architecture diagram showing XLM-RoBERTa encoder with 6 prediction heads\\
\textbf{Option 2:} Performance comparison across tasks (bar chart: accuracy/F1 by classification task)\\
\textbf{Option 3:} Confidence score distribution showing high-confidence vs. flagged-for-review predictions\\
\textbf{Option 4:} Conceptual workflow diagram: source DBs → model → classified policies → human review loop
\vspace{100pt}}}
\caption{[Replace with actual figure]. Suggested visualizations: (1) multi-task model architecture, (2) classification performance by task, (3) confidence score distributions, or (4) end-to-end pipeline workflow.}
\label{fig:placeholder}
\end{figure}

\end{document}
